\documentclass[12pt,a4paper,twoside]{article}

%\usepackage[utf8]{inputenc}
\usepackage[T1]{fontenc}
\usepackage[english]{babel}
\usepackage{lmodern}
\usepackage{tikz}

\usetikzlibrary{positioning,calc,arrows}

\usepackage{listings}

\usepackage{amsmath}
\usepackage[labelformat=simple]{subfig}
\usepackage[font={sf},margin=10pt,labelfont=bf]{caption}
\usepackage{booktabs}
\usepackage[colorlinks=false]{hyperref}

\usepackage{fancyhdr}
\setlength{\headheight}{15pt}
\usepackage{todonotes, minted}

\clubpenalty = 10000
\widowpenalty = 10000 \displaywidowpenalty = 10000

\begin{document}

\def\sectionautorefname{Section}
\def\subsectionautorefname{Subsection}

\newcommand\ie{i.e.\ }
\newcommand\eg{e.g.\ }

  \begin{titlepage}
    \begin{tikzpicture}[remember picture,overlay]
      % Seitenrahmen zeichnen.
      \draw[semithick,rounded corners=0.5cm]
        ($(current page.north west) + ( 1cm,-1cm)$) --
        ($(current page.north east) + (-1cm,-1cm)$) --
        ($(current page.south east) + (-1cm, 1.5cm)$);

      \draw[semithick,rounded corners=0.5cm]
        ($(current page.south east) + (-1cm, 1.5cm)$) --
        ($(current page.south west) + ( 1cm, 1.5cm)$) --
        ($(current page.north west) + ( 1cm,-1cm)$);

      % Logo einbinden.
      \node[anchor=north west] (logo) at ($(current page.north west) + (1.75cm,-1.5cm)$)
      {
        \includegraphics[width=4cm]{KITLogo}
      };

      % Institut / Lehrstuhl.
      \node[anchor=east] at ($(current page.east |- logo.east) + (-1.75cm,0cm)$)
      {
        \begin{minipage}{5.2cm}
          \begin{flushleft}
            \footnotesize{}Institut für Programmstrukturen und Datenorganisation (IPD) \\
            \vspace{6pt}
            Lehrstuhl Prof.\ Dr.-Ing.\ Snelting
          \end{flushleft}
        \end{minipage}
      };

      \node (title) at ($(current page.center |- logo.south) + (0cm, -4cm)$)
      {
        % Korrekter Zeilenabstand etc. durch Minipage.
        \begin{minipage}{12cm}
          \begin{center}
            \huge\textbf{Simple Verification of Rust Programs via Functional Purification}
          \end{center}
        \end{minipage}
      };

      \node[below=1.75cm of title.south]   (prename)  { Masterarbeit von };
      \node[below=0.75cm of prename.south] (name)     { \Large{}\textbf{Sebastian Ullrich} };
      \node[below=1cm    of name.south]    (postname) { an der Fakultät für Informatik };

      \node[below=3cm    of name.south]    (bildchen) { \includegraphics[width=0.9\textwidth]{logo.png}
                                                      };

      \node[below=2cm of bildchen.south] (table)
      {
        \begin{tabular}{ll}
          \textbf{Erstgutachter:}           & Prof.\ Dr.-Ing.\ Gregor Snelting \\[5pt]
          \textbf{Zweitgutachter:}          & ??? \\[5pt]
        \end{tabular}
      };

      \node[below=3.5cm of table.south] (time)
      {
        \begin{tabular}{ll}
        \textbf{Bearbeitungszeit:} & 4. Juli 2013 -- 29. Oktober 2013
        \end{tabular}
      };

      % Fußzeile, unten zentriert.
      \node[anchor=south] (footnote) at ($(current page.center |- current page.south) + (0cm, 0.65cm)$)
      {
        \tiny{}KIT -- Universität des Landes Baden-Württemberg und nationales Forschungszentrum in der Helmholtz-Gemeinschaft
        \hspace{0.5cm}
        \Large{}\textbf{www.kit.edu}
      };
    \end{tikzpicture}
  \end{titlepage}

% sane default for proof documents
\parindent 0pt\parskip 0.5ex

\tikzset{every node/.style={transform shape},auto,block/.style={align=center,rectangle,draw,minimum height=20pt,minimum width=30pt},>=triangle 60}
%\pagenumbering{Roman}
\pagestyle{empty}
\renewcommand{\abstractname}{Einfache Verifikation von Rust-Programmen}
\begin{abstract}
  Imperative Programmiersprachen sind in der modernen Softwareentwicklung
  allgegenwärtig, stellen aber ein Hindernis für formale Softwareverifikation
  dar durch ihre Verwendung von veränderbaren Variablen und Objekten. Programme
  in diesen Sprachen können normalerweise nicht direkt auf die unveränderliche
  Welt von Logik und Mathematik zurückgeführt werden, sondern müssen in eine
  explizit modellierte Semantik der jeweiligen Sprache eingebettet werden. Diese
  Indirektion stellt ein Problem für die Benutzung von interaktiven
  Theorembeweisern dar, da sie die Entwicklung von neuen Werkzeugen, Taktiken
  und Logiken für diese ``innere'' Sprache bedingt.

  Die vorliegende Arbeit stellt einen Compiler von der imperativen
  Programmiersprache Rust in die pur funktionale Sprache des Theorembeweisers
  Lean vor, der nicht nur generell das erste Werkzeug zur Verifikation von Rust-Programmen
  darstellt, sondern diese insbesondere auch
  mithilfe der von Lean bereitgestellten Standardwerkzeugen und -logik
  ermöglicht. Diese Transformation ist nur möglich durch spezielle Eigenschaften
  von allen validen Rust-Programmen, die die Veränderbarkeit von Werten auf
  begrenzte Geltungsbereiche einschränken und statisch durch Rusts Typsystem
  garantiert werden. Die Arbeit demonstriert den Einsatz des Compilers anhand
  der Verifikation von Realbeispielen und zeigt die Erweiterbarkeit des Projekts
  über reine Verifikation hinaus am Beispiel von asymptotischer Laufzeitanalyse auf.
\end{abstract}
\renewcommand{\abstractname}{Abstract}

\newpage

\begin{abstract}
  Imperative programming, and aliasing in particular, represents a major
  obstacle in formally reasoning about everyday code. By utilizing restrictions
  the imperative programming language Rust imposes on mutable aliasing, we
  present a scheme for transforming a substantial part of the Rust language into the
  purely functional language of the Lean theorem prover. We use this scheme to
  verify the correctness of real-world examples of Rust code without the need
  for special semantics or logics. We furthermore
  show the extensibility of our transformation by incorporating an analysis of
  asymptotic runtimes.
\end{abstract}
\tableofcontents

\cleardoublepage
\pagestyle{fancy}
\fancyhf{}
\fancyhead[LE,RO]{\thepage}
\fancyhead[RE,LO]{\textit\leftmark}
\pagenumbering{arabic}

\section{Introduction}

Imperative programming languages are ubiquitous in today's software development,
making them prime targets for formal reasoning. Unfortunately, their semantics
differ from those of mathematics and logic -- the languages of formal methods -- in some
significant details, starting with the very concept of ``variables''. The problem
of mutability is only exacerbated for languages that allow references to
\emph{alias}, or point to the same memory location, enabling non-local mutation.

The standard way of verifying programs in such languages with the help of an
interactive theorem prover is to explicitly model the semantics of the language
in the language of the theorem prover, then translate
the program to this representation (a ``deep'' embedding) and finally prove the
correctness of its formalized behavior. This general approach is very flexible
and allows for the verification of meta programs such as program transformations.
The downside of the approach is that the theorem prover's tools and tactics may
not be directly applicable to the embedded language, defeating many amenities of
modern theorem provers.
Alternatively, programs can be ``shallowly'' embedded by directly translating them
into terms in the theorem prover's language without the use of an explicit inner semantics.
This simplifies many semantic details such as the identification and
substitution of bound variables, but it is harder to accomplish the more the semantics
of the source language differs from the theorem prover's own semantics.

Regardless of the type of embedding,
an explicit heap that references can point into must generally be modeled and
passed around in order to deal with the aliasing problem. References in this model may be as simple as
indices into a uniform heap, but various logics such as separation logic~\cite{reynolds2002separation} have been developed to work on a more abstract representation and to express
aliasing-free sets of references.

Languages with more restricted forms of aliasing exist, however.
Rust~\cite{matsakis2014rust}, a new, imperative systems programming language,
imposes on mutable references the restriction of never being aliased by any
other reference, mutable or immutable. This restriction eliminates the
possibility of data races and other common bugs created by the presence of
mutable sharing such as iterator invalidation. It furthermore enables more
aggressive optimizations.

While the full Rust language also provides raw pointers, which are not bound by
the aliasing restriction, and other ``unsafe'' operations, a
memory model for Rust (informal or formal) has yet to be proposed. We therefore focus on the ``safe''
subset of Rust that has no unsolved semantic details.

We utilize safe Rust's aliasing restriction to design a monadic shallow embedding of a
substantial subset of Rust
into the purely functional language of the Lean~\cite{de2015lean} theorem prover, without the need
for any heap-like indirection. This allows us to
reason about unannotated, real-world Rust code in mostly the same manner one would
reason about native Lean definitions. The monadic approach gives us further
flexibility in modeling additional effects such as function runtime.

We first discuss the simpler cases of the
translation, notably excluding mutable references, in \autoref{sec:trans}. We
show their application by giving a formal verification of Rust's
\verb![T]::binary_search! method in \autoref{sec:binary_search}.
\autoref{sec:mutref} discusses the translation of most usages of mutable
references, which is used in \autoref{sec:fixedbitset} for a verification of the
\texttt{fixedbitset} crate.
\newpage
\section{Related Work}
\label{sec:related}

While this thesis presents the first verification tool for Rust programs, tools
for many other imperative languages have been developed before.

The Why3 project~\cite{bobot2011why3} is notable for its generality. It provides
an imperative ML-like language \emph{WhyML} together with a verification
condition generator that can interface with a multitude of both automatic and
interactive theorem provers. While WhyML supports advanced language features such
as type polymorphism and exceptions, it does not support higher-order functions,
which are ubiquitous in Rust code.
%Specification annotations of WhyML programs are written in the logic language \emph{Why3} and thus cannot 
WhyML provides a reference type \texttt{ref} that can point to a fresh cell on
the heap and is statically checked not to alias with other \texttt{ref}
instances, but cannot point into some existing datum like Rust references can.
For example, the first of the following two WhyML functions fails to type check
because the array elements are not known to be alias-free, while the second one
will return a reference to a \emph{copy} of \verb!a[i]!.

\begin{minted}{sml}
let get_mut (a : array (ref int)) (i : int) : ref int = a[i]
let get_mut (a : array int) (i : int) : ref int = ref a[i]
\end{minted}

In contrast, Rust can provide a perfectly safe function with this functionality.

\begin{minted}{rust}
fn get_mut<T>(slice: &mut [T], index: usize) -> &mut T
\end{minted}

WhyML is also being used as an intermediate language for the verification of
programs in Ada~\cite{guitton2011hi}, C~\cite{cuoq2012frama} and Java~\cite{filliatre2007krakatoa}.
For the latter two languages, aliasing is reintroduced by way of an explicit heap.

The remarkable SeL4 project~\cite{klein2009sel4} delivers a full formal verification of an operating
system microkernel by way of multiple levels of program verification and
refinement steps. The C code that produces the final kernel binary is verified
by embedding it into the theorem prover
Isabelle/HOL~\cite{nipkow2002isabelle}, using a deep embedding for statements
and a shallow one for expressions. The C memory model used allows type-unsafe
operations by use of a byte-size heap, but as with Why3, higher-order functions are
not supported. The AutoCorres~\cite{greenaway2012bridging, greenaway2014don}
tool then transforms this representation into a shallow monadic embedding,
dealing with the `uninteresting complexities of C'~\cite{greenaway2014don} on the
way. The result is an abstracted representation that is quite similar to ours
(and in fact inspired it in part, as we shall note below), but does not go the
last mile of completely eliminating the heap where possible. Thus the user still
has to worry and reason about (the absence of) aliasing manually or through a
nested logic such as separation logic. Without explicit
no-alias annotations, the semantics of C would allow eliminating the heap in far fewer places than those
of Rust in any case.

It should be noted that our work, like most verification projects based on
either embedding or code extraction, relies on both
an unverified compiler and an unverified embedding tool, effectively making both
part of the trusted computing base. SeL4 is a notable exception in this,
providing (at lower optimization levels) a direct equivalence proof~\cite{sewell2013translation} between the
produced kernel binary and the verified embedded code, thus completely removing
the original C code from the trusted computing base.

While not an imperative language, the purely functional, total Cogent language~\cite{o2016refinement}
uses linear types in the style of Wadler~\cite{wadler1990linear} for safe
manual memory management, much like Rust. The language is designed both to be
easily verifiable (by building on AutoCorres) and to
compile down to efficient C code. As we shall see in \autoref{sec:rust}, the
biggest differences between Wadler-style purely functional linear languages and Rust are the
existence of mutable references as well as sophisticated interprocedural
reference tracking in the latter. For example, the aforementioned \rust{get_mut}
function can only be expressed as a higher-order function in Cogent, even in the immutable case.
\newpage

\section{Background}
\subsection{Rust}
\subsection{Lean}

\section{Basic Transformation}
\label{sec:trans}

\section{Case Study: Verification of \texttt{std::[T]::binary\_search}}
\label{sec:binary_search}

\section{Transformation of Mutable References}
\label{sec:mutref}

\section{Case Study: Verification of \texttt{fixedbitset}}
\label{sec:fixedbitset}

\section{Conclusion}

\newpage
\bibliographystyle{abbrv}
\bibliography{bib}

\cleardoublepage
\pagestyle{empty}
\section*{Erklärung}

  \vspace{20mm}
  Hiermit erkläre ich, Sebastian Andreas Ullrich, dass ich die vorliegende Masterarbeit selbst\-ständig
verfasst habe und keine anderen als die angegebenen Quellen und Hilfsmittel
benutzt habe, die wörtlich oder inhaltlich übernommenen Stellen als solche kenntlich gemacht und
die Satzung des KIT zur Sicherung guter wissenschaftlicher Praxis beachtet habe.
  \vspace{20mm}
  \begin{tabbing}
  \rule{4cm}{.4pt}\hspace{1cm} \= \rule{7cm}{.4pt} \\
 Ort, Datum \> Unterschrift
  \end{tabbing}

\end{document}
